\chapter{Introduction}
\label{chap:introduction}

\section{Motivation}
\label{sec:motivation}

When developers work with AI coding agents, the requirements of a project are
rarely expressed in a single instruction. Instead, the developer's intent
evolves across a long and unstructured prompt history containing initial
requests, follow-up questions, corrections, removals, replacements, and changes
of direction. Although this history records how the project developed, it does
not provide a clear description of what the project should be in its current
state. Earlier requirements may no longer apply, later prompts may replace
previous values, and some messages may concern temporary implementation
decisions rather than the intended product.

To make the project easier to continue or transfer, developers can create a
\emph{master prompt}. In this thesis, a master prompt is a single,
self-contained specification of the user's currently valid intent. It should
describe what must be built without requiring another developer or AI agent to
read the entire prompt history. It should keep active requirements, reflect
later changes, and leave out requirements that were removed or replaced.

Recently, AI coding agents can build much of a software project from a clear set
of prompts. This makes the prompts used to describe the project an important
part of the project itself. Work on spec-driven development also shows that a
written specification can be kept and used to guide the generated code
\citep{piskala2026specdriven}. In the same way, a master prompt can be treated
as a useful project artifact rather than a temporary summary. It can be saved,
updated, and kept together with the code so that the project can be continued
in another AI session, by another coding agent, or by another developer.

However, creating a faithful master prompt is not only a summarization task.
The agent must separate what the user actually requested from details that
appeared during implementation. In practical coding, an AI agent may add architectural
patterns, filenames, technologies, validation rules, or other choices that the
user never asked for. This problem is closely related to \emph{extrinsic hallucination}, which
\citet{ji2023survey} describe as generated content that is not supported by
the source. In master-prompt consolidation, these unsupported additions may
make the result appear more complete while reducing its faithfulness to the
user's original requests.

The main concern is therefore not only whether an AI agent can produce a
clear and compact document, but whether every requirement in that document is
traceable to the prompt history. A faithful master prompt should preserve all
currently valid requirements, correctly reflect removals and replacements, and
exclude both agent-added details and outdated requirements. Although carefully
written consolidation instructions may improve the result, it remains unclear
how strongly instruction strictness affects traceability and whether stricter
prompting continues to improve the result for longer and more realistic
development histories.

A second challenge arises when the master prompt is treated not only as the
output of a one-time consolidation process, but as a persistent project
artifact. After handoff, a later session or another contributor may introduce a
new instruction that contradicts a requirement already established in the
master prompt. Such contradictions may use different wording that may unintentionally
modify, replace, or even remove already existing requirements without recognizing the contradiction.
Moreover, they may appear in different parts of the master prompt rather than next to each
other. They are therefore semantic conflicts rather than text-level merge conflicts.
If the agent silently chooses one value, combines incompatible requirements,
or overwrites the existing requirement, the master prompt can no longer serve
as a reliable representation of the project's intent.

Together, these problems motivate a spec-guided approach that keeps the current requirements
and their sources explicit, and checks each new instruction against the
existing specification before updating it.

\section{Problem Statement}
\label{sec:problem-statement}

This thesis studies two connected problems in AI-assisted master prompt
construction:

\begin{itemize}[leftmargin=*]
\item \textbf{Non-traceable consolidation.}
This problem occurs when an agent consolidates the full prompt history. The
resulting master prompt may contain details that the user never requested,
such as technologies, filenames, design choices, or added rules. It may also
retain old requirements after they were removed, undone, or replaced. At the
same time, some valid requirements may be omitted. These errors make the
master prompt less faithful even when its wording is clear and its overall
meaning appears similar to the original history.

\item \textbf{Semantic conflict in a persistent master prompt.}
This problem occurs when a new prompt is added after the master prompt has
already been created. The new instruction may contradict an existing
requirement, but the agent may fail to report the conflict. Instead, it may
silently overwrite the old value, combine two incompatible values, or keep
both requirements in different parts of the document. In each case, the
master prompt no longer provides one consistent description of the project.
\end{itemize}

To address these problems, this thesis evaluates several single-pass
consolidation instructions and proposes a \emph{spec-guided} approach. Rather
than consolidating the raw prompt history only at the end of a session, the
proposed approach maintains an explicit specification of the current user
intent under a governance instruction. Each new prompt is checked against the
existing requirements before the specification is updated. Potential conflicts
are surfaced together with their source requirements instead of being silently
resolved.

\section{Research Questions}
\label{sec:research-questions}

\begin{description}[leftmargin=*,style=nextline]
\item[RQ1]
To what extent do AI-generated master prompts contain non-traceable content,
including agent-added requirements and outdated requirements that should
have been removed or replaced? How does consolidation instruction strictness
(\emph{basic}, \emph{loose}, and \emph{strict}) affect this problem?

\item[RQ2]
Does a spec-guided consolidation approach produce a more traceable master
prompt than single-pass consolidation of the raw prompt history using basic,
loose, or strict instructions?

\item[RQ3]
When a later instruction semantically conflicts with a settled requirement,
how can a spec-guided approach surface the conflict, trace it to both
requirements, and avoid resolving it without user input?
\end{description}

\section{Contributions}
\label{sec:contributions}

This thesis makes the following contributions:

\begin{itemize}[leftmargin=*]

\item A requirement-level evaluation method for master-prompt traceability.
Instead of comparing whole documents, the method splits each master prompt
into atomic requirements and uses entailment-based scoring with AlignScore
\citep{zha2023alignscore} to determine whether each requirement is supported
by the human-consolidated ground truth. Paired examples are also used to show
why word similarity alone is not a reliable measure of traceability. Applied
to 24 master-prompt candidates produced by three agents, four consolidation
approaches, and two projects, the method answers RQ1 by showing that stricter
single-pass instructions reduce non-traceable content, but the improvement
stops before reaching the ground truth
(\cref{chap:experiment-traceability}).

\item A consolidation method called \texttt{spec\_guided}, which maintains a
living specification under an explicit governance instruction. Instead of
rebuilding the master prompt from the full prompt history only at the end of
the session, the method checks each new prompt and updates the specification
throughout the session. This method answers RQ2. Compared with each agent's
own single-pass baseline, it improves traceability in four of the six
agent--project pairs. Manual analysis of the other two pairs identifies
trade-offs caused by making requirement decisions repeatedly during the
session rather than once at the end
(\cref{chap:experiment-traceability}).

\item A project-handoff method for identifying semantic conflicts in a settled
master prompt. The method uses a separate governance instruction that checks
each later prompt against the existing specification, preserves settled
requirements when the intended change is unclear, and records both conflicting
requirements for the user to resolve. The method is demonstrated using five
constructed conflicts and two compatible controls across the simple and hard
projects (\cref{chap:experiment-conflict}).

\end{itemize}
